\par\medskip
\noindent c) Hay que tomar la imagen de este homomorfismo por el
operador de conexión $\partial$ de (2.3.7), asociado a la sucesión
exacta (2.3.6). Ahora bien,
$nQ\simeq Q/{}_nQ$ es también manifiestamente isomorfo a $M_0/M'_0$,
de modo que tenemos un homomorfismo de sucesiones exactas

%%%%%%%%%%  scan idx 257  |  folio 246  %%%%%%%%%%

\[
\begin{tikzcd}[column sep=large,row sep=large]
 0\arrow[r]&M'_0\arrow[r]\arrow[d]&M_0\arrow[r]\arrow[d]
   &M_0/M'_0\arrow[r]\arrow[d,"\wr"']&0\\
 0\arrow[r]&{}_nQ\arrow[r]&Q\arrow[r]&nQ\arrow[r]&0&{,}
\end{tikzcd}
\]
donde la última flecha vertical es un isomorfismo. Esto muestra que la
sección buscada de $\uExt^1(nQ,G)$ también puede calcularse como la
imagen de $\varphi'$ de (2.3.14) por el operador de conexión
\[
 \uHom(M'_0,G)\xrightarrow{\ \partial\ }
 \uExt^1(M_0/M'_0,G),
\]
o, equivalentemente, como la sección del segundo término definida por
el elemento de $\Ext^1$ global que es la imagen de $\varphi'$ por el
operador de conexión
\[
 \Hom(M'_0,G)\longrightarrow
 \Ext^1(M_0/M'_0,G).
\]
La construcción usual del operador de conexión da la receta siguiente.
Prolónguese el compuesto (2.3.13)
\[
 M'_0\xrightarrow{\ \varphi'\ }G\longrightarrow C^0
\]
a un homomorfismo
\[
 \psi:M_0\longrightarrow C^0,
\]
lo cual es posible porque $C^0$ es inyectivo, y considérese el único
homomorfismo
\begin{equation}\tag{2.3.15}
 \lambda:M_0/M'_0\longrightarrow C^1
\end{equation}
que hace conmutativo el diagrama
\[
\begin{tikzcd}[column sep=large,row sep=large]
 0\arrow[r]&M'_0\arrow[r]\arrow[d,"\varphi'"']
   &M_0\arrow[r]\arrow[d,"\psi"']
   &M_0/M'_0\arrow[r]\arrow[d,"\lambda"']&0\\
 0\arrow[r]&G\arrow[r]&C^0\arrow[r]&C^1&&{.}
\end{tikzcd}
\]

%%%%%%%%%%  scan idx 258  |  folio 247  %%%%%%%%%%

Es claro que (2.3.15) toma valores en los $1$-cociclos, es decir, que
es un $1$-cociclo, y que este cociclo representa el elemento buscado de
\[
 \Ext^1(nQ,G)=\Ext^1(M_0/M'_0,G).
\]

\medskip
\noindent d) Queda por describir explícitamente la sección de
$\uExt^1(nQ,G)$ obtenida de $\xi$ y $x$ mediante la flecha (2.3.5)
del diagrama (2.3.10). Para ello podemos olvidar que la sección $x$ de
$P$ está anulada por $n$. En términos del levantamiento elegido $x_0$
de $x$ a una sección de $L_0$, la descripción es la siguiente. Este
levantamiento define un homomorfismo de aumentación
\[
 u:M_\bullet\longrightarrow L_\bullet\otimes M_\bullet,
 \qquad
 y_0\longmapsto x_0\otimes y_0,
 \quad
 y_1\longmapsto x_0\otimes y_1.
\]
La sección de $\uExt^1(Q,G)$ obtenida de (2.3.1) es, por tanto, la que
se obtiene del elemento $\xi$ de
$\Ext^1(L_\bullet\otimes M_\bullet,G)$ componiendo con el
homomorfismo precedente. Está representada, pues, por el homomorfismo
de grado uno
\[
\begin{tikzcd}[column sep=large,row sep=huge]
 0\arrow[r]&M_1\arrow[r,"d_1"]
   \arrow[d,"{y_1\mapsto f_1(x_0\otimes y_1)}"']
   \arrow[rd,phantom,"-" description]
   &M_0\arrow[r]
   \arrow[d,"{y_0\mapsto f_{\sisi{1}{0}}(x_0\otimes y_0)}"]&0\\
 0\arrow[r]&C^0\arrow[r,"d^0"]&C^1&\ldots&{.}
\end{tikzcd}
\]
Para restringir esta clase a $nQ\hookrightarrow Q$, utilicemos el
isomorfismo $nQ\simeq M_0/M'_0$, que da una resolución plana
\[
 0\longrightarrow M'_0\longrightarrow M_0\longrightarrow0
\]
de $nQ$. La inclusión $nQ\hookrightarrow Q$ se levanta al
homomorfismo de resoluciones siguiente:
\[
\begin{tikzcd}[column sep=large,row sep=huge]
 0\arrow[r]&M'_0\arrow[r]
   \arrow[d,"{y'_0\mapsto[ny'_0]}"']
   \arrow[rd,phantom,"+" description]
   &M_0\arrow[r]\arrow[d,"{y_0\mapsto ny_0}"]&0\\
 0\arrow[r]&M_1\arrow[r]&M_0\arrow[r]&0&{.}
\end{tikzcd}
\]

%%%%%%%%%%  scan idx 259  |  folio 248  %%%%%%%%%%

Por composición, la sección buscada de $\uExt^1(nQ,G)$ está, pues,
asociada al elemento de $\Ext^1(nQ,G)$ representado por el
homomorfismo de complejos de grado uno
\[
\begin{tikzcd}[column sep=large,row sep=large]
 0\arrow[r,dash]&M'_0\arrow[r,dash,"d_1"]\arrow[d,"g_1"']
   &M_0\arrow[r,dash]\arrow[d,"g_0"]&0&{}\\
 0\arrow[r,dash]&C^0\arrow[r,dash]&C^1\arrow[r,dash]&\dots&{,}
\end{tikzcd}
\]
donde
\[
 g_1(y'_0)=f_1\bigl(x_0\otimes[ny'_0]\bigr),
 \qquad
 g_0(y_0)=n f_0(x_0\otimes y_0).
\]

\medskip
\noindent e) Utilizando ahora que $x$ es una sección de $\,{}_nP$,
mostremos que el elemento precedente de $\Ext^1(nQ,G)$ es el opuesto
del elemento definido por (2.3.15). Para todo homomorfismo
\[
 h:M_0\longrightarrow C^0,
\]
el elemento de $\Ext^1$ considerado en d) no cambia si $g_1,g_0$ se
reemplazan, respectivamente, por
\[
 g'_1=g_1+h d_1,
 \qquad
 g'_0=g_0-d^0h.
\]
Disponemos del homomorfismo $\psi:M_0\longrightarrow C^0$ construido
en c), que por definición satisface
\[
 \psi(y'_0)=
 f_1\bigl([nx_0]\otimes y'_0-x_0\otimes[ny'_0]\bigr)
\]
para toda sección $y'_0$ de $M'_0$. Equivalentemente,
\[
 f_1\bigl(x_0\otimes[ny'_0]\bigr)
 \bigl(=g_1(y'_0)\bigr)
 =f_1\bigl([nx_0]\otimes y'_0\bigr)-\psi(y'_0).
\]
Definamos $h:M_0\longrightarrow C^0$ por

%%%%%%%%%%  scan idx 260  |  folio 249  %%%%%%%%%%

\[
 h(y_0)=-f_1\bigl([nx_0]\otimes y_0\bigr)+\psi(y_0).
\]
Por construcción se tiene entonces
\[
 g'_1=0,
 \qquad
 g'_0(y_0)=n f_0(x_0\otimes y_0)
 +d^0f_1\bigl([nx_0]\otimes y_0\bigr)-d^0\psi(y_0).
\]
Como $d^0f_1=-f_0d_1$, tenemos
\[
 d^0f_1\bigl([nx_0]\otimes y_0\bigr)
 =-f_0(nx_0\otimes y_0),
\]
y por tanto
\[
 g'_1=0,
 \qquad
 g'_0=-d^0\psi.
\]
Así, el elemento considerado es el definido por el homomorfismo
\[
 M_0/M'_0\simeq nQ\longrightarrow C^1
\]
inducido por $-d^0\psi$ al pasar al cociente. Obtenemos, por tanto, el
opuesto del elemento definido por (2.3.15).

\medskip
\noindent\underline{Observaciones 2.3.16.}
La demostración precedente da en realidad una compatibilidad algo más
precisa que la indicada en 2.3.11, obtenida eligiendo una sección $x$
de $\,{}_nP$. Esta sección define entonces un homomorfismo
\[
 {}_nQ\longrightarrow G,
 \qquad
 y\longmapsto\varphi_\xi(x,y),
\]
y, por el homomorfismo de conexión asociado a (2.3.6), un elemento de
$\Ext^1(nQ,G)$. Por otra parte, olvidando que $x$ está anulado por
$n$, consideramos el homomorfismo $\mathbb Z_T\longrightarrow P$
definido por $\sisi{X}{x}$. Su imagen inversa es una biextensión de
$(\mathbb Z_T,Q)$ por $G$, cuya clase puede interpretarse como un
elemento de $\Ext^1(Q,G)$ y, por tanto, induce un elemento de
$\Ext^1(nQ,G)$. Con estas convenciones, los dos elementos obtenidos
en $\Ext^1(nQ,G)$ son \underline{opuestos}.
